% ============================================================
% ============================================================
%  OVERALL ASSESSMENT
% ============================================================
% ============================================================

\section[Overview]{\textbf{NACA0012 Overall Assessment}}

% ============================================================
% Overall Assessment
% ============================================================

\begin{frame}{\small Overall Assessment}
\scriptsize

\begin{block}{\scriptsize Overall Assessment}
\begin{itemize}
    \item \textbf{Aerodynamics:} Predictions are well clustered; $C_p$ agrees well with experiment, while $C_D$ shows the greatest grid sensitivity.

    \item \textbf{Droplet Impingement:} Limited grid sensitivity is observed, with 7-bin results closely reproducing the 15-bin reference.

    \item \textbf{HTC \& Thermodynamics:} Local $HTC$ exhibits substantial code-to-code variability, influenced by turbulence modeling, roughness, and computation methodology; surface temperature and freezing fraction remain highly variable regardless of grouping.

    \item \textbf{Ice Accretion:} Ice mass is relatively consistent, but the experimental AE3933 mass reduction is substantially underpredicted; final ice geometry remains strongly code-dependent.
\end{itemize}
\end{block}

\vspace{0.15cm}

\begin{alertblock}{\scriptsize Overall Takeaway}
Aerodynamics, impingement, and ice mass are comparatively consistent across submissions, while larger variability emerges in thermodynamics and final ice geometry. The AE3932--AE3933 change in horn orientation is captured, but the experimental reduction in ice mass remains substantially underpredicted.
\end{alertblock}

\end{frame}